\documentclass[11pt]{article}
\usepackage[margin=1in]{geometry}
\usepackage{lineno}
\usepackage{booktabs}
\usepackage{enumitem}
\usepackage{xcolor}
\usepackage{titling}
\usepackage{titlesec}
\usepackage{tabularx}
\usepackage{multirow}
\usepackage{float}
\usepackage{graphicx}
\usepackage[font=small,labelfont=bf,labelsep=period,
  justification=raggedright,singlelinecheck=false]{caption}
\usepackage[colorlinks=true,linkcolor=black,urlcolor=blue!60!black,citecolor=black]{hyperref}

% Inline references as clickable [N] links (no bibliography section)
\newcommand{\citelink}[2]{\href{#1}{\small[#2]}}

\renewcommand\linenumberfont{\normalfont\scriptsize\sffamily\color[rgb]{0.5,0.5,0.5}}
\leftlinenumbers
\linenumbers

\titleformat{\section}{\large\bfseries}{\thesection.}{0.5em}{}
\titlespacing{\section}{0pt}{2.0em}{0.6em}
\titlespacing{\subsection}{0pt}{1.8em}{0.5em}

\setlength{\droptitle}{-4em}
\pretitle{\begin{center}\large\bfseries}
\posttitle{\end{center}\vspace{-1em}}
\preauthor{\begin{center}\normalsize}
\postauthor{\end{center}\vspace{-1em}}
\predate{\begin{center}\normalsize}
\postdate{\end{center}\vspace{-0.5em}}

\title{Results}
\author{Diego Toribio\\[0.3em]{\normalsize March 10, 2026}}
\date{\vspace{-2.5em}}

\begin{document}
\maketitle
\thispagestyle{empty}

\noindent{\small This document is a working snapshot of experimental
results. DQN isolation study and Rainbow baseline are complete across
six games, seed 42. Rainbow+SPR runs are pending.}

\section{Introduction}

Atari-100K \citelink{https://arxiv.org/abs/1903.00374}{1} is a
data-efficiency benchmark for reinforcement learning. The agent gets
just 100K environment steps per game -- roughly two hours of real-time
play. That is two orders of magnitude less experience than the standard
protocol used to train the original DQN
\citelink{https://www.nature.com/articles/nature14236}{2}, so the
question becomes: can an agent learn anything useful with so little
data?

\medskip
An \textit{environment step} is one decision point: the agent sees a
game screen, picks an action, and receives a reward. Under the standard
frame skip of 4, each decision is held for 4 consecutive video frames,
so 100K steps correspond to 400K rendered frames.

\medskip
SPR \citelink{https://arxiv.org/abs/2007.05929}{3} adds a
self-supervised objective that trains the agent to predict its own
future internal states, encouraging richer representations even when
reward is scarce. However, SPR was only evaluated on top of Rainbow
\citelink{https://arxiv.org/abs/1710.02298}{4} -- a complex agent with
many components -- so it is unclear whether the representation learning
is what helped or whether Rainbow did the heavy lifting. This leads to
our central question:

\begin{quote}
When does self-supervised representation learning improve data-efficient
reinforcement learning, and what determines whether the learned
representations are useful?
\end{quote}

\medskip
We first isolate SPR on vanilla DQN, using \textit{data augmentation}
(small random shifts to each game frame) as a second factor. By testing
each method alone and in combination across six games, we can separate
their individual and combined effects. We then add Rainbow+SPR runs on
the same games to provide the contrasting case where SPR is known to
produce large improvements. Comparing representations across these two
base agents -- one simple, one capable -- reveals what the base agent
provides that makes self-supervised representations useful.

\section{Methods}

\subsection{Game Selection}

We selected six games from the Atari-100K benchmark to cover a range of
reward structures and visual challenges, ensuring that any patterns in
how SPR and augmentation interact are not artifacts of a single game
type. The selection was guided by published Rainbow+SPR results in
Schwarzer et al.\ \citelink{https://arxiv.org/abs/2007.05929}{3}, which
show diverse interaction patterns across these games.

\begin{table}[H]
\centering
\renewcommand{\arraystretch}{1.3}
\setlength{\tabcolsep}{10pt}
\begin{tabular}{l l}
\toprule
\textbf{Game} & \textbf{Challenge} \\
\midrule
Crazy Climber & Continuous vertical scrolling; dense reward from climbing \\
Boxing        & Two-player combat; dense reward from landing punches \\
Kangaroo      & Platformer with enemies; moderate reward spacing \\
Road Runner   & Fast horizontal scrolling; frequent item collection \\
Frostbite     & Platform assembly over ice; delayed reward structure \\
Up N Down     & Vertical driving; dense reward from jumping on cars \\
\bottomrule
\end{tabular}
\caption{The six games and their primary challenges. Games were chosen
to span dense to moderate reward structures and a variety of visual
dynamics.}
\end{table}

\subsection{Experimental Conditions}

We start with a vanilla DQN baseline and progressively add augmentation
and SPR. Each condition will be run across six games with three seeds
(72 runs total); results reported here use seed 42.

\begin{table}[H]
\centering
\renewcommand{\arraystretch}{1.3}
\setlength{\tabcolsep}{12pt}
\begin{tabular}{l l}
\toprule
\textbf{Condition} & \textbf{What changes} \\
\midrule
Baseline  & Vanilla DQN, no modifications \\
+ Aug     & Add random-shift data augmentation \\
+ SPR     & Add self-predictive representation learning \\
+ Both    & Augmentation and SPR combined \\
\bottomrule
\end{tabular}
\caption{The four experimental conditions. Each builds on the same DQN
architecture and hyperparameters.}
\end{table}

\subsection{Training Setup}

All conditions use the Nature CNN architecture
\citelink{https://www.nature.com/articles/nature14236}{2} as the shared
encoder and are trained with the same configuration. Augmentation and
SPR add their own components on top of this base. Results report mean
raw episode return ($\pm$ std) across three seeds.

\begin{table}[H]
\centering
\renewcommand{\arraystretch}{1.2}
\setlength{\tabcolsep}{12pt}
\begin{tabular}{l l}
\toprule
\textbf{Parameter} & \textbf{Value} \\
\midrule
Optimizer       & RMSProp (lr = 0.00025) \\
Batch size      & 32 \\
Exploration     & $\epsilon$: 1.0 $\to$ 0.1 \\
Replay buffer   & 100K capacity, 2K warmup \\
Target network  & Hard copy every 2K steps \\
\bottomrule
\end{tabular}
\caption{Training configuration shared across all four DQN conditions.}
\end{table}

\noindent The agent explores by choosing random actions with probability
$\epsilon$, which decays from 1.0 (fully random) to 0.1 over the first
half of training. The replay buffer stores past transitions so the agent
can retrain on earlier experience; training begins after 2K transitions
have been collected. The target network is a frozen copy of the main
network used to compute stable learning targets, refreshed every 2K
steps.

\subsection{Data Augmentation}

The augmentation condition uses DrQ-style random shifts
\citelink{https://arxiv.org/abs/2004.14990}{5}: each 84$\times$84
grayscale frame is zero-padded to 92$\times$92 and randomly cropped
back to 84$\times$84, shifting the image by up to 4 pixels in any
direction. The shift is applied independently to each frame in a
sampled training batch, so the agent sees slightly different views of
the same transition on each update. No augmentation is applied during
evaluation.

\subsection{Self-Predictive Representations (SPR)}

SPR adds a self-supervised auxiliary loss that trains the encoder to
predict its own future representations. A transition model (two
convolutional layers with action conditioning) predicts the next latent
state, and a cosine similarity loss compares the prediction against
targets produced by an exponential moving average (EMA) copy of the
encoder (momentum 0.99). Predictions are made over 5 future steps.
Projection and prediction heads (MLPs) sit on top of the encoder to
avoid representational collapse. The SPR loss is added to the standard
TD loss with weight $\lambda = 2.0$. Encoder dropout (0.5) is used for
regularization when augmentation is not present.

\subsection{Rainbow Comparison}

To test whether SPR needs a capable base agent, we run the same six
games with Rainbow DQN -- a stronger agent that combines six
improvements over vanilla DQN: distributional value estimates (C51,
51 atoms), double Q-learning, dueling network streams, noisy
exploration, 3-step returns, and prioritized experience replay.

\begin{table}[H]
\centering
\renewcommand{\arraystretch}{1.3}
\setlength{\tabcolsep}{12pt}
\begin{tabular}{l l}
\toprule
\textbf{Condition} & \textbf{What changes} \\
\midrule
Rainbow       & Rainbow DQN, no SPR \\
Rainbow + SPR & Add self-predictive representation learning \\
\bottomrule
\end{tabular}
\caption{Rainbow conditions. Both share the same architecture and
hyperparameters; SPR is the only difference.}
\end{table}

\noindent Rainbow uses Adam (lr = $6.25 \times 10^{-5}$) instead of
RMSProp, updates the target network every 8K steps, and requires 20K
warmup transitions for prioritized replay to stabilize. All other
settings (100K environment steps, 6 games) match the DQN conditions.

\section{Results}

All results report mean raw episode return ($\pm$ std) over 30
evaluation episodes at the final checkpoint (400K frames), seed 42.

\begin{table}[H]
\centering
\renewcommand{\arraystretch}{1.4}
\setlength{\tabcolsep}{8pt}
\begin{tabular}{l rrrr r}
\toprule
\textbf{Game}
  & \textbf{Baseline}
  & \textbf{+ Aug}
  & \textbf{+ SPR}
  & \textbf{+ Both}
  & \textbf{Rainbow} \\
\midrule
Crazy Climber  & $3540 \pm 754$   & $110 \pm 40$      & $3400 \pm 0$      & $3.3 \pm 18$     & $9137 \pm 3857$  \\
Road Runner    & $1440 \pm 859$   & $3.3 \pm 18$      & $2290 \pm 1238$   & $1033 \pm 539$   & $180 \pm 464$    \\
Boxing         & $-3.0 \pm 5.3$   & $-18.2 \pm 4.0$   & $-18.2 \pm 11.7$  & $-39.6 \pm 10$   & $-3.0 \pm 6.5$   \\
Kangaroo       & $26.7 \pm 68$    & $460 \pm 128$     & $13.3 \pm 49.9$   & $320 \pm 98$     & $33.3 \pm 90.7$  \\
Frostbite      & $147 \pm 45$     & $0.0 \pm 0.0$     & $161 \pm 31$      & $173 \pm 15$     & $43 \pm 18$      \\
Up N Down      & $141 \pm 155$    & $981 \pm 671$     & $901 \pm 682$     & $491 \pm 467$    & $1794 \pm 637$   \\
\midrule
Mean           & $982$  & $256$  & $1125$  & $330$  & $1864$ \\
\bottomrule
\end{tabular}
\caption{Summary of all conditions. Mean raw episode return ($\pm$
std) over 30 eval episodes at 400K frames, seed 42. The first four
columns are DQN variants; Rainbow is a separate base agent.}
\label{tab:summary}
\end{table}

\noindent The learning curves below reveal instability that the summary
table obscures. Most conditions swing between high and near-zero
performance across checkpoints, making final-checkpoint comparisons
fragile. Several patterns stand out:

\subsection{Augmentation}

Augmentation is strongly game-dependent. Kangaroo ($460$ vs $27$) and
Up N Down ($981$ vs $141$) improve substantially, while Crazy Climber,
Road Runner, and Frostbite collapse to near-zero returns. Boxing worsens
slightly. The pattern suggests that random shifts help in games where
spatial invariance aids generalization but hurt when precise spatial
information (e.g.\ climber position, road layout) is critical for the
policy.

\begin{figure}[H]
\centering
\includegraphics[width=\textwidth]{../output/plots/aug_helps.pdf}
\caption{Games where augmentation helps. Kangaroo shows augmentation
building steadily while SPR peaks early then collapses. Up N Down shows
both augmentation and SPR improving over baseline, with SPR climbing
late in training.}
\label{fig:aug-helps}
\end{figure}

\subsection{SPR}

SPR's effect on vanilla DQN is modest. Up N Down shows a strong lift
($901$ vs $141$, roughly $6\times$) and Road Runner improves ($2290$ vs
$1440$). The remaining four games are flat or slightly worse: Crazy
Climber holds steady, Frostbite gains marginally, and both Boxing and
Kangaroo decline. This contrasts with the large improvements reported in
the original SPR paper, which used Rainbow as the base agent rather than
vanilla DQN.

\begin{figure}[H]
\centering
\includegraphics[width=\textwidth]{../output/plots/aug_hurts_spr_holds.pdf}
\caption{Games where augmentation hurts but SPR holds. Crazy Climber
shows baseline and SPR both learning while augmentation flatlines.
Road Runner shows SPR emerging as the strongest condition late in
training while augmentation collapses early.}
\label{fig:aug-hurts}
\end{figure}

\subsection{Combined}

The combined condition (+Both) uses the same EMA momentum (0.99) and
dropout (0.5) as the SPR-only condition, with augmentation as the sole
addition. Augmentation does not rescue SPR on this base agent -- the
games where augmentation helps (Kangaroo, Up N Down) see reduced gains
when combined with SPR, and the games where augmentation hurts (Crazy
Climber, Road Runner) still collapse.

\begin{figure}[H]
\centering
\includegraphics[width=\textwidth]{../output/plots/minimal_learning.pdf}
\caption{Games with minimal learning across conditions. Boxing shows all
conditions declining, with the combined condition performing worst.
Frostbite shows flat performance except for a brief spike in the
combined condition that does not sustain.}
\label{fig:minimal-learning}
\end{figure}

\subsection{Rainbow Baseline}

Rainbow is not uniformly better under the 100K budget. Crazy Climber
($9137$ vs $3540$, roughly $3\times$) and Up N Down ($1794$ vs $141$,
$12\times$) see large gains. Both are dense reward games where
Rainbow's components (distributional targets, prioritized replay,
n-step returns) appear to calibrate within 100K steps.

\begin{figure}[H]
\centering
\includegraphics[width=\textwidth]{../output/plots/rainbow_helps.pdf}
\caption{Games where Rainbow improves over DQN. Crazy Climber shows
Rainbow consistently above baseline, peaking higher before both
collapse at the final checkpoint. Up N Down shows Rainbow pulling
ahead in the second half of training.}
\label{fig:rainbow-helps}
\end{figure}

\noindent Road Runner ($180$ vs $1440$) and Frostbite ($43$ vs $147$)
regress substantially. Both games performed reasonably under vanilla
DQN, so the regression is not a floor effect -- the added complexity
hurts more than it helps with so little data.

\begin{figure}[H]
\centering
\includegraphics[width=\textwidth]{../output/plots/rainbow_hurts.pdf}
\caption{Games where Rainbow regresses from DQN. Road Runner shows
DQN peaking mid-training while Rainbow stays near zero throughout.
Frostbite shows both agents oscillating at similar levels, with
Rainbow dipping to zero at 200K before recovering.}
\label{fig:rainbow-hurts}
\end{figure}

\noindent Boxing and Kangaroo are effectively unchanged -- both agents
score near zero regardless of the base agent. These games appear too
difficult for either architecture within 100K steps.

\begin{figure}[H]
\centering
\includegraphics[width=\textwidth]{../output/plots/rainbow_flat.pdf}
\caption{Games where neither agent learns. Kangaroo shows both agents
near zero with high variance. Boxing shows Rainbow recovering to
near-zero while DQN declines steadily.}
\label{fig:rainbow-flat}
\end{figure}

\subsection{Summary}

Three patterns emerge from the single-seed results. First, SPR's
effect on vanilla DQN is modest and game-dependent -- it helps on two
of six games and is flat or slightly harmful on the rest. This
contrasts with the large improvements reported when SPR is paired with
Rainbow. Second, augmentation is similarly game-dependent but in
different directions: it helps games where spatial invariance aids
generalization and collapses games where precise spatial information is
critical. Third, Rainbow improves the same two games where SPR and
augmentation also showed their strongest effects (Crazy Climber, Up N
Down), suggesting that dense reward games with relatively simple visual
dynamics are most responsive to additional capability under the 100K
budget. The key open question is whether adding SPR to Rainbow produces
the large improvements seen in published results, and if so, what
Rainbow provides at the representation level that vanilla DQN does not.

\section{Next Steps}

All results so far use a single seed (42). Three areas of work remain
before the thesis experiments are complete.

\subsection{Multi-Seed Evaluation}

Each condition will be run with seeds 42, 13, and 23 to produce
confidence intervals. The current single-seed results show high
variance between adjacent checkpoints, so multi-seed aggregation is
necessary to distinguish real effects from noise. Learning curves will
be re-plotted with shaded confidence bands across seeds.

\subsection{Rainbow + SPR}

The Rainbow baseline establishes that a more capable agent improves on
some games but regresses on others. The next step adds SPR to Rainbow
on the same six games. Comparing DQN+SPR (modest effects) against
Rainbow+SPR (large effects in published results) directly tests whether
the base agent's capability determines whether self-supervised
representations are useful.

\subsection{Interpretability Analysis}

The performance comparison alone cannot explain why SPR helps Rainbow
but not vanilla DQN. Four analyses are planned to investigate the
representation-level differences:

\medskip
\begin{enumerate}[leftmargin=*, itemsep=0.6em]
  \item \textbf{Linear probes.} Train classifiers on frozen encoder
    outputs across all conditions (DQN, DQN+SPR, Rainbow, Rainbow+SPR)
    to measure what game-relevant information the representations
    encode.
  \item \textbf{Latent space visualization.} t-SNE or PCA projections
    of encoder outputs across conditions, testing whether SPR produces
    structured representations only when paired with a capable base
    agent.
  \item \textbf{Transition model predictions.} Compare SPR's predicted
    next-states to actuals for DQN+SPR and Rainbow+SPR. If prediction
    quality differs, the bottleneck is in the representations; if
    not, it is downstream.
  \item \textbf{Component attribution.} Correlate Rainbow training
    metrics (distributional loss, importance sampling weights,
    priority entropy) with probe accuracy to identify which
    components drive representation quality, then run targeted
    ablations on the strongest candidates.
\end{enumerate}

\clearpage

\appendix
\section*{Appendix}
\addcontentsline{toc}{section}{Appendix}

\section{Known Issue: Reward Clipping in Evaluation}
\label{app:clipping}

An early three-game validation run (Crazy Climber, Kung Fu Master,
Boxing) reported returns in \textbf{clipped reward space} ($\{-1, 0,
+1\}$ per step) because the evaluation environment inherited the
training reward clipper. Published baselines
\citelink{https://arxiv.org/abs/2007.05929}{3}%
\citelink{https://arxiv.org/abs/2004.14990}{5}
report \textbf{raw Atari scores}. The two scales are not comparable --
for example, the random agent scores 10,780 raw on Crazy Climber while
the clipped-reward agent reported 153. This was fixed for all runs in
Table~\ref{tab:summary} by disabling reward clipping in the evaluation
environment.

\section{Early Baseline Observations}
\label{app:observations}

Detailed per-game analysis from the initial three-game validation run
(clipped rewards):

\medskip
\begin{enumerate}[leftmargin=*, itemsep=0.6em]
  \item \textbf{Crazy Climber shows clear learning.} Monotonic
    improvement through 320K with variance-driven dips. Loss converged
    from 0.87 to 0.055 (moving average). Viable baseline anchor for the
    comparative study.

  \item \textbf{Kung Fu Master shows no learning.} Eval never exceeds
    3.7 (clipped). 458 episodes in 400K frames indicates the agent
    dies quickly ($\sim$870 frames per episode). Loss converged normally
    (2.49$\to$0.024), suggesting the network fits Q-values but they do
    not translate to useful behavior within the 100K budget.

  \item \textbf{Boxing learned a losing policy.} Eval oscillates
    wildly ($-0.2$ to $-42.5$), with the agent actively losing by large
    margins in late training. Only 56 episodes in 400K frames
    ($\sim$7100 frames each) meant the agent saw fewer than 28 episodes
    before epsilon reached 0.1 at 200K frames.

  \item \textbf{Epsilon decay may be too aggressive for long-episode
    games.} By 200K frames, Crazy Climber had seen $\sim$90 episodes
    (sufficient), Kung Fu Master $\sim$230 (sufficient but short-lived),
    and Boxing only $\sim$28 (insufficient for an 18-action game).
\end{enumerate}

\section{Game Selection Reference}
\label{app:games}

The 26 Atari-100K benchmark games grouped by reward density -- how
often the agent receives a non-zero score during play. In dense games
the agent scores on most actions; in sparse games, long stretches pass
with no reward. Of the 26 games, 7 are dense, 12 moderate, and
4 sparse. Categorization based on Schwarzer et al.\
\citelink{https://arxiv.org/abs/2007.05929}{3}, Table 4.

\vspace{0.3em}
\begin{table}[H]
\centering
\renewcommand{\arraystretch}{1.15}
\setlength{\tabcolsep}{14pt}
\begin{tabular}{l l l}
\toprule
\textbf{Dense} & \textbf{Moderate} & \textbf{Sparse} \\
\midrule
Road Runner     & Seaquest    & Pong        \\
Crazy Climber   & Frostbite   & Freeway     \\
Kung Fu Master  & Hero        & Private Eye \\
Up N Down       & Alien       & Bank Heist  \\
Krull           & Asterix     &             \\
Demon Attack    & Ms.\ Pac-Man &            \\
Chopper Command & BattleZone  &             \\
                & Kangaroo    &             \\
                & Breakout    &             \\
                & Jamesbond   &             \\
                & Boxing      &             \\
                & Gopher      &             \\
\bottomrule
\end{tabular}
\caption{Atari-100K benchmark games grouped by how frequently the agent
receives non-zero reward. Dense games yield reward on most steps; sparse
games have long stretches of zero reward.}
\end{table}

\end{document}
